%-------------------------
% Resume in Latex
% Author : Jake Gutierrez
% Based off of: https://github.com/sb2nov/resume
% License : MIT
% Build : tectonic -X compile template.tex
%         (XeTeX; fonts via fontspec, Unicode text in the PDF stays ATS-parsable)
%------------------------

\documentclass[letterpaper,11pt]{article}

\usepackage{latexsym}
\usepackage[empty]{fullpage}
\usepackage{titlesec}
\usepackage{marvosym}
\usepackage[usenames,dvipsnames]{color}
\usepackage{verbatim}
\usepackage{enumitem}
\usepackage[hidelinks]{hyperref}
\usepackage{fancyhdr}
\usepackage[english]{babel}
\usepackage{tabularx}
\usepackage{fontawesome5}
\usepackage{multicol}
\usepackage{fontspec}
\usepackage{kotex} % xetexko under XeTeX
\setlength{\multicolsep}{-3.0pt}
\setlength{\columnsep}{-1pt}


%----------FONT OPTIONS----------
% sans-serif
% \usepackage[sfdefault]{FiraSans}
% \usepackage[sfdefault]{roboto}
% \usepackage[sfdefault]{noto-sans}
% \usepackage[default]{sourcesanspro}

% serif
% \usepackage{CormorantGaramond}
\setmainfont{XCharter}[ % OpenType extension of Bitstream Charter
  Extension      = .otf,
  UprightFont    = *-Roman,
  BoldFont       = *-Bold,
  ItalicFont     = *-Italic,
  BoldItalicFont = *-BoldItalic,
]
\xetexkofontregime{latin} % parentheses/punctuation in the Latin font, not the Hangul one
\frenchspacing            % no extra space after sentence punctuation


\pagestyle{fancy}
\fancyhf{} % clear all header and footer fields
\fancyfoot{}
\renewcommand{\headrulewidth}{0pt}
\renewcommand{\footrulewidth}{0pt}

% Adjust margins
\addtolength{\oddsidemargin}{-0.6in}
\addtolength{\evensidemargin}{-0.5in}
\addtolength{\textwidth}{1.19in}
\addtolength{\topmargin}{-.7in}
\addtolength{\textheight}{1.4in}

\urlstyle{same}

\raggedbottom
\raggedright
\setlength{\tabcolsep}{0in}

% Sections formatting
\titleformat{\section}{
\vspace{-4pt}\scshape\raggedright\large\bfseries
}{}{0em}{}[\color{black}\titlerule \vspace{-5pt}]

%-------------------------
% Custom commands
\newenvironment{ieeereflist}{
\begin{list}{}{
\setlength{\leftmargin}{2.0em}
\setlength{\itemindent}{0pt}
\setlength{\labelwidth}{1.5em}
\setlength{\labelsep}{0.5em}
\setlength{\listparindent}{0pt}
\setlength{\parsep}{0pt}
\setlength{\topsep}{0pt}
\setlength{\partopsep}{0pt}
}
}{
\end{list}
}
\newcommand{\ieeerefitem}[2]{\item[\small\textbf{[#1]}] \small#2}

\newcommand{\resumeItem}[1]{
\item\small{
{#1 \vspace{-2pt}}
}
}

\newcommand{\classesList}[4]{
\item\small{
{#1 #2 #3 #4 \vspace{-2pt}}
}
}

\newcommand{\resumeSubheading}[4]{
\vspace{-2pt}\item
\begin{tabular*}{1.0\textwidth}[t]{l@{\extracolsep{\fill}}r}
\textbf{#1} & \textbf{\small #2} \\
\textit{\small#3} & \textit{\small #4} \\
\end{tabular*}\vspace{-7pt}
}

\newcommand{\resumeSubheadingInline}[5]{
\vspace{-2pt}\item
\begin{tabular*}{1.0\textwidth}[t]{l@{\extracolsep{\fill}}r}
\textbf{#1}, {#2} & \textbf{\small #3} \\
\textit{\small#4} & \textit{\small #5} \\
\end{tabular*}\vspace{-7pt}
}

\newcommand{\resumeSubheadingInlineTwo}[5]{
\vspace{-2pt}\item
\begin{tabular*}{1.0\textwidth}[t]{l@{\extracolsep{\fill}}r}
\textbf{#1} $|$ {#2} & \textbf{\small #3} \\
\textit{\small#4} & \textit{\small #5} \\
\end{tabular*}\vspace{-7pt}
}

\newcommand{\resumeSubSubheading}[2]{
\item
\begin{tabular*}{0.97\textwidth}{l@{\extracolsep{\fill}}r}
\textit{\small#1} & \textit{\small #2} \\
\end{tabular*}\vspace{-7pt}
}

\newcommand{\resumeProjectHeading}[2]{
\item
\begin{tabular*}{1.001\textwidth}{l@{\extracolsep{\fill}}r}
\small#1 & \textbf{\small #2}\\
\end{tabular*}\vspace{-7pt}
}

\newcommand{\resumeSubItem}[1]{\resumeItem{#1}\vspace{-4pt}}

\renewcommand\labelitemi{$\vcenter{\hbox{\tiny$\bullet$}}$}
\renewcommand\labelitemii{$\vcenter{\hbox{\tiny$\bullet$}}$}

\newcommand{\resumeSubHeadingListStart}{\begin{itemize}[leftmargin=0.0in, label={}]}
\newcommand{\resumeSubHeadingListEnd}{\end{itemize}}
\newcommand{\resumeItemListStart}{\begin{itemize}[leftmargin=0.2in]}
\newcommand{\resumeItemListEnd}{\end{itemize}\vspace{-5pt}}

\newcommand{\resumeProjectSubheading}[3]{
\vspace{1pt}\item[]
\hspace*{-0.5\labelwidth}%
\begin{tabular*}{\dimexpr\linewidth+0.5\labelwidth\relax}[t]{l@{\extracolsep{\fill}}r}
\textbf{\small #1} {#2} & {\small #3} \\
\end{tabular*}\vspace{-5pt}
}

%-------------------------------------------
%%%%%%  RESUME STARTS HERE  %%%%%%%%%%%%%%%%%%%%%%%%%%%%


\begin{document}

\begin{center}
{\Huge \scshape David Hyunyoo Jang} \\ \vspace{1pt}
\small 
\href{mailto:hwjang00@snu.ac.kr}{\raisebox{-0.2\height}\faEnvelope\ {hwjang00@snu.ac.kr}} ~ 
\href{https://linkedin.com/in/daylight00/}{\raisebox{-0.2\height}\faLinkedin\ {daylight00}}  ~
\href{https://github.com/daylight-00}{\raisebox{-0.2\height}\faGithub\ {daylight-00}} ~ $|$ ~ US Citizen
\vspace{-8pt}
\end{center}


%-----------EDUCATION-----------
\section{Education}
\resumeSubHeadingListStart
\resumeSubheading
{Seoul National University}{Mar 2022 -- Feb 2028 (Expected)}
{PharmD}{Seoul, South Korea}
\resumeSubheading
{Korea University}{Mar 2019 -- Aug 2021}
{Coursework in Materials Science and Engineering (no degree)}{Seoul, South Korea}
\resumeSubHeadingListEnd


%-----------EXPERIENCE-----------
\section{Research Experience}
\resumeSubHeadingListStart
\resumeSubheadingInline
{AI Protein Design Lab}{Korea Advanced Institute of Science and Technology (KAIST)}{Jun 2025 -- Aug 2025}
{Undergraduate Researcher (advisor: Gyu Rie Lee)}{Daejeon, South Korea}
\resumeItemListStart
\resumeProjectSubheading{De Novo Binder Design for Phenytoin and Carbamazepine Sensors}{}{}
\resumeItem{Conducted computational de novo binder-design campaigns for phenytoin and carbamazepine sensors in collaboration with Kirill Alexandrov's lab; generated and screened candidates using diffusion-based backbone sampling, LigandMPNN, and Rosetta, then delivered shortlisted designs to the collaborators.}
\resumeItem{Consolidated the binder-design workflow into SLURM-parallelized stages with parquet-based candidate-level metric tracking; integrated AF3/Boltz-2 refolding and PLACER analysis, and implemented symmetry-aware, linker-agnostic ligand RMSD to evaluate binding poses after protein alignment. [\href{https://github.com/daylight-00/ligand_binder_design}{\color{RoyalBlue}{code}}]}
% Note: pool-level numbers (21->275 screened, -25.4->-31.6) and the interface ddG trend over cycles are not used: the former are means of final pools with different cutoffs and backbone counts, the latter is the optimized objective with no orthogonal check. The ~32K count is 63 independent trajectories x 32 cycles x 16 candidates = 32,256 FastRelax runs (24 no-linker + 39 linker starting structures, CBZ/PHT; recomputed from the handoff parquets and the run scripts); it is workload, not an optimization result. Baseline check: the 2022 original_script samples 1 sequence per structure in intermediate cycles and runs bsite repack + MinMover(bb,chi,jump) with FastRelax only in the final cycle, so its structures already update every cycle; do not write "beyond fixed-backbone". Ours: 16 candidates per cycle, each FastRelaxed (side-chain packing on), one state chosen by ddg (interface Ddg filter, ascending) carried forward. Intent vs result (user, 2026-11-07): GRL's script aimed at sequence diversity; the added selection was meant to move the sampling space in a favorable direction, but no orthogonal check was ever run and the user later concluded it is not optimization (design_cycle was started for that). So state the purpose as intent ("to explore") and never claim an effect; do not lead with "extended an existing script" and do not write "the lab's ..." anywhere in this block (user, 2026-11-07). The bullet below is the external reviewer's third proposal (2026-11-07), as written; the ~32K count is no longer stated (63 x 16 x 32 = 32,256 is still derivable). Pool means were also cross-checked against the raw stats JSON of the linker run (39 trajectories, identical to the parquet). Verified basis: over the 63 trajectories the pool MEANS (16 candidates per cycle) fall from cycle 1 to cycle 32 for ddg (-30.24 -> -35.13, lower in 62/63 trajectories) and for the LigandMPNN score (0.799 -> 0.674, 60/63); res_totalscore and cms move the same way but are not stated. In-loop metrics only: no control arm and no per-cycle refolding or ipTM, so no causal or binder-quality claim. Data caveat: fast_relax_parallel appends results in as_completed order while the selection step and the parse notebook index them by sequence index, so in fastrelax.parquet the per-row pairing of (mpnn, seq, path) with (ddg, cms, res_totalscore) is unreliable (in the linker run the entry at the true selected index carried the selected ddg in only 106 of 1248 cycles, about chance). Set-level pool means, best-of-16 ddg/cms/res_totalscore (the recycled structure) and the recycled structure itself are unaffected. Never state per-candidate mpnn-vs-ddg results or the mpnn of the selected candidate.
\resumeItem{Developed metric-guided LigandMPNN--FastRelax recycling to steer sequence sampling beyond diversification, selecting the lowest-Rosetta-ddG structure from 16 relaxed candidates per cycle; across 63 trajectories, observed decreases in mean candidate-pool Rosetta ddG ($-30.24 \to -35.13$) and LigandMPNN scores ($0.799 \to 0.674$) from cycle 1 to 32. [\href{https://github.com/daylight-00/ligandMPNN_FR}{\color{RoyalBlue}{code}}]}
\resumeProjectSubheading{RFdiffusion2 Atom-wise rASA Conditioning}{}{}
\resumeItem{Extended RFdiffusion2 inference to support atom-wise relative solvent accessibility (rASA) conditioning, enabling control over which ligand atoms are buried or exposed to guide binder orientation; implemented ligand atom-name mapping, preserved global-rASA compatibility, and demonstrated the capability on a trimethyllysine (M3L) target. [\href{https://github.com/RosettaCommons/RFdiffusion2/pull/35}{\color{RoyalBlue}{PR}}]}
\resumeItemListEnd

\resumeSubheadingInline
{Laboratory of Computational Drug Discovery}{Seoul National University}{Jul 2023 -- Aug 2025}
{Undergraduate Researcher (advisor: Juyong Lee)}{Seoul, South Korea}
\resumeItemListStart
\resumeProjectSubheading{PREpiBind --- pMHC-II binding: representation benchmark \& web server}{}{}
\resumeItem{Led PREpiBind, a controlled comparison of 10 protein representations under a shared joint-attention architecture for pMHC-II binding across qualitative, mass-spectrometry, and IC50 datasets; protein language models (PLMs) generally led qualitative/MS benchmarks, while NetMHCIIpan-4.3 performed best on thresholded IC50.}
\resumeItem{Designed leakage-controlled splits and held-out evaluations (leave-one-molecule-out and mouse H2), showing that PLM advantages narrowed under distribution shift; reimplemented DeepNeo for baseline comparison.}
\resumeItem{Showed that ESMC 300M matched ESMC 600M and ESM3 Small/Medium within seed variation (ROC-AUC 0.919--0.921); released code, datasets, trained checkpoints, and a public inference server with custom-sequence support and Colab/CLI demos. [\href{https://github.com/daylight-00/PREpiBind}{\color{RoyalBlue}{code}}; \href{https://lcdd.snu.ac.kr/prepibind/}{\color{RoyalBlue}{server}}]}
\resumeProjectSubheading{LLMs on the Korean Pharmacist Licensing Exam}{}{}
\resumeItem{Led an end-to-end bilingual benchmark of 27 proprietary and open-source LLMs on six years (2019--2024) of the Korean Pharmacist Licensing Examination, developing the exam extraction, translation, local/API inference, and official-scoring pipeline with human-percentile comparisons.}
\resumeItem{Found that seven models passed all six exams in both languages, with the top model ranking among the top 12\% of human examinees; characterized logarithmic model size--performance trends and assessed potential pretraining contamination using temporal, cross-language, and subject-level evidence. [\href{https://github.com/daylight-00/llms-pass-kple}{\color{RoyalBlue}{code}}]}
\resumeProjectSubheading{CB1 receptor small-molecule virtual screening}{}{}
\resumeItem{Reproduced an end-to-end structure-based virtual screen of the Enamine Hit Locator Library (460K compounds) against the CB1 receptor on an 8-GPU cluster ($\sim$12 h) --- self-compiling AutoDock-GPU from source for multi-GPU execution (the prebuilt binary used 1 of 8 GPUs) --- and automating batch docking, result parsing, and ranking.}
\resumeItemListEnd

\resumeSubheadingInline
{Structure \& Knowledge Injection into Machine Learning Lab}{Seoul National University}{May 2024 -- Sep 2024}
{Undergraduate Researcher (advisor: Jay-Yoon Lee)}{Seoul, South Korea}
\resumeItemListStart
\resumeProjectSubheading{Semi-supervised data augmentation for RoseTTAFold2NA}{\small\textit{collaboration with Minkyung Baek's lab}}{}
\resumeItem{Implemented the training pipeline for a semi-supervised joint-loss objective that augments RF2NA with unlabeled self-distilled structures --- built the mixed labeled/unlabeled data loader, per-batch loss routing, and a distributed sampler with tunable unlabeled fraction, consolidating scattered prototype scripts into one parameterized training entry point.}
\resumeItem{Designed and ran a controlled low-label study --- PDB-only vs distilled-augmentation vs joint-loss semi-supervision across 10/30/50\% label-density regimes --- to test whether unlabeled data improves RF2NA learning under sparse supervision.}
\resumeItemListEnd

\resumeSubheadingInline
{Systems Pharmacology Lab}{Seoul National University}{Jun 2023 -- Jul 2023}
{Undergraduate Internship (advisor: Minsoo Noh)}{Seoul, South Korea}
\resumeItemListStart
\resumeItem{Performed HaCaT cell culture and PGE$_2$ competitive ELISA to assess anti-inflammatory effects of tryptophol; conducted dose--response analysis and benchmarking against NSAIDs.}
\resumeItemListEnd
\resumeSubHeadingListEnd


%-----------PUBLICATIONS AND PREPRINTS-----------
\section{Publications \& Preprints}
\begin{ieeereflist}
\ieeerefitem{1}{PREpiBind: Protein Representation-integrated Epitope--MHC Class II Binding Prediction\\
\textbf{David Hyunyoo Jang}, Dongwoo Kim, Untaek Hwang, Byungho Park, Yoonjoo Choi, Juyong Lee\\
Submitted; preprint: \textit{bioRxiv} (2026), \href{https://doi.org/10.64898/2026.09.15.751749}{\color{RoyalBlue}{doi:10.64898/2026.09.15.751749}}}
\ieeerefitem{2}{Proprietary and Open-Source Large Language Models on the Korean Pharmacist Licensing Examination: A Comparative Benchmarking Study\\
\textbf{David Hyunyoo Jang}, Juyong Lee\\
Under review; preprint: \textit{medRxiv} (2025), \href{https://doi.org/10.1101/2025.04.15.25325584}{\color{RoyalBlue}{doi:10.1101/2025.04.15.25325584}}}
\end{ieeereflist}


%-----------SELECTED RESEARCH PROJECTS-----------
% DRAFT: entries only; bullets are written after each project's final review.
\section{Selected Research Projects}
\resumeSubHeadingListStart
\resumeSubheading{foldspace}{2026 -- present}{Self-initiated; evaluating pretrained folding representations as generative priors}{}
\resumeSubheading{design\_cycle}{2026 -- present}{Self-initiated; follow-up to LigandMPNN--FastRelax recycling}{}
\resumeSubHeadingListEnd


%-----------RESEARCH SOFTWARE AND OPEN SOURCE-----------
\section{Research Software \& Open-Source Contributions}
% top-level list: use the nested-list spacing of the other sections' bullets
\begin{itemize}[leftmargin=0.2in, itemsep=2pt, parsep=2pt]
\resumeItem{\textbf{ESMFold2--AtomWorks} --- integrated ESMFold2 with AtomWorks data structures and an optional Foundry training path; exact feature parity on reference fixtures and exact output parity under deterministic GPU kernels. [\href{https://github.com/daylight-00/esmfold2-atomworks}{\color{RoyalBlue}{code}}]}
\resumeItem{\textbf{RosettaCommons Foundry} --- merged LigandMPNN inference improvements: per-design and ligand-interface confidence reporting, and parser overrides for AtomArray inputs. [\href{https://github.com/RosettaCommons/foundry/pull/306}{\color{RoyalBlue}{\#306}}; \href{https://github.com/RosettaCommons/foundry/pull/389}{\color{RoyalBlue}{\#389}}]}
\resumeItem{\textbf{Upstream fixes} --- AtomWorks insertion-code residue parsing, proposed and incorporated upstream in \#116 [\href{https://github.com/RosettaCommons/atomworks/pull/107}{\color{RoyalBlue}{\#107}}]; ESMFold2 multi-component CCD ligand output, the same change now present upstream [\href{https://github.com/Biohub/esm/pull/416}{\color{RoyalBlue}{\#416}}].}
\resumeItem{\textbf{patinae} (Rust/WGPU molecular viewer) --- fixed renderer failures on low-limit GPUs and Windows. \textit{Merged} [\href{https://github.com/zmactep/patinae/pull/28}{\color{RoyalBlue}{\#28}}; \href{https://github.com/zmactep/patinae/pull/29}{\color{RoyalBlue}{\#29}}]}
\resumeItemListEnd % = \end{itemize}\vspace{-5pt}


%-----------PRESENTATIONS-----------
\section{Presentations}
\begin{ieeereflist}
\ieeerefitem{1}{PREpiBind: Protein Representation-integrated Epitope--MHC Class II Binding Prediction\\
KIDDS 2025 Summer Workshop (Korean In silico bioDesign and Discovery Society), Poster}
\ieeerefitem{2}{DeepXeno: Structure-Informed Prediction of Xenotransplantation Immunogenicity \textit{(project later developed into PREpiBind)}\\
2024 SNU Student-Directed Education Undergraduate Research Program, Talk (Outstanding Research Presentation)}
\end{ieeereflist}


%-----------AWARDS AND FUNDING-----------
\section{Awards \& Funding}
\resumeSubHeadingListStart
\resumeSubheadingInline{Research Excellence Award}{2024 SNU Student-Directed Education Undergraduate Research Program}{2025}{President of Seoul National University}{}
\resumeItemListStart
\resumeItem{Awarded to two projects per fields (ten total) for excellence in research execution and outcomes.}
\resumeItemListEnd
\resumeSubheading{2024 SNU Student-Directed Education Undergraduate Research Program}{2024}{Faculty of Liberal Education, Seoul National University}{}
\resumeItemListStart
\resumeItem{Selected as one of 36 projects to receive funding for undergraduate research. \textit{(KRW 7M)}}
\resumeItemListEnd
\resumeSubHeadingListEnd


%-----------TEACHING AND LEADERSHIP-----------
\section{Teaching \& Leadership}
\resumeSubHeadingListStart
\resumeSubheadingInlineTwo{SNU Alpha}{Deep Learning Study Club}{Spring 2022 -- Winter 2024}{Executive Member, Co-Founder}{College of Pharmacy, Seoul National University}
\resumeItemListStart
\resumeItem{Co-founded and helped organize a peer-led deep learning study club for pharmacy students.}
\resumeItem{Led weekly study sessions on Natural Language Processing (Fall 2023) and PyTorch (Spring 2024), teaching deep learning concepts and applications in computational drug discovery.}
\resumeItemListEnd
\resumeSubheading
{Military Service}{Aug 2025 -- Feb 2027 (Expected)}
{Team Leader, Sergeant}{III Corps, Republic of Korea Army}
\resumeItemListStart
\resumeItem{Led a 10-person air-defense team, coordinating personnel and overseeing the execution of assigned duties.}
\resumeItem{Operated and maintained air-defense radar systems in a frontline sector along the Korean DMZ.}
\resumeItemListEnd
\resumeSubHeadingListEnd


%-----------CLINICAL PHARMACY TRAINING-----------
\section{Clinical Pharmacy Training}
\resumeSubHeadingListStart
\resumeSubheading
{Seoul National University Hospital}{Dec 2024 -- Feb 2025}
{Hospital Pharmacy Practicum}{Seoul, South Korea}
\resumeItemListStart
\resumeItem{Completed rotations across 12 specialty departments, participating in medication dispensing, drug therapy case reviews, and patient counseling under pharmacist supervision.}
\resumeItemListEnd
\resumeSubheading
{E-Blue Pharmacy}{Feb 2025 -- Mar 2025}
{Community Pharmacy Practicum}{Suwon, South Korea}
\resumeItemListStart
\resumeItem{Reviewed prescriptions for potential medication-related problems and participated in medication therapy reviews and counseling on prescription and over-the-counter medications under pharmacist supervision.}
\resumeItemListEnd
\resumeSubHeadingListEnd


%-----------TECHNICAL SKILLS-----------
\section{Technical Skills}
\begin{itemize}[leftmargin=0.15in, label={}]
\small{\item{
\textbf{Programming \& Data}: Python, PyTorch, Numpy, Bash, pandas, matplotlib, Streamlit, Nginx \\
\textbf{Modeling \& Docking}: RDKit, PyMOL, Chimera, AutoDock, PyRosetta, RFdiffusion (AA), LigandMPNN, AlphaFold3, ESMFold2/ESMC, AtomWorks/Foundry, TMol, Biotite \\
\textbf{Systems \& DevOps}: Linux, Git, Slurm, conda/mamba, HuggingFace \\
\textbf{Experimental Techniques}: cell culture (HaCaT), buffer/reagent prep, competitive ELISA, serial dilutions, plate reader
}}
\end{itemize}


\end{document}
